\documentclass[a5paper,10pt]{article}

% https://github.com/InDevRus/filippov-solutions

% Нумерация страниц
\pagenumbering{gobble}

% Настройка полей
\usepackage{geometry}
\geometry{left=1cm}
\geometry{right=1cm}
\geometry{top=1cm}
\geometry{bottom=1.5cm}

% Вставка таблицы
\usepackage{array}
\usepackage{tabularx}

% Инструменты выравнивания формул
\usepackage{amsmath}
\usepackage{mathtools}

% Диакритика в математике
\usepackage{amsfonts}

% Вычисления размеров на ходу
\usepackage{calc}

% Удобные записи производных
\usepackage[italic=true]{derivative}

% Настройки сносок
\usepackage{enotez}

% Длинные знаки векторов
\usepackage{anyfontsize}
\usepackage[b]{esvect}

% Вставка картинок
\usepackage{graphicx}

% Вставка красной строки
\usepackage{indentfirst}
\setlength{\parindent}{1em}

% Условия в командах
\usepackage{xifthen}

% Луа-скрипты
\usepackage{luacode}

% Настройка команд отдельно в математическом и текстовом режимах
\usepackage{mathcommand}

% Для повторения LaTeX кода
\usepackage{multido}

% Конфигурация отступов формул
\delimitershortfall-1sp
\usepackage{mleftright}
\mleftright

% Растягивание объектов
\usepackage{relsize}

% Обтекание текстом
\usepackage{wrapfig}
\usepackage{subcaption}
\usepackage{floatflt}

% Поддержка ввода кириллицы
\usepackage[english, russian]{babel}

% Настройка корневой позиции для компилятора
\usepackage{import}
\providecommand*{\ProjectRootPath}{..}

\import{\ProjectRootPath/preambles/macros/}{theorems}

\import{\ProjectRootPath/preambles/fonts}{font_setup}

% Знак градуса
\usepackage{siunitx}

\import{\ProjectRootPath/preambles/macros/}{operators}
\import{\ProjectRootPath/preambles/macros/}{redefinitions}

\import{\ProjectRootPath/preambles/fonts}{letters_setup}
\import{\ProjectRootPath/preambles/fonts/adjustments}{subsuperscripts}

\import{\ProjectRootPath/preambles/custom}{frames}
\import{\ProjectRootPath/preambles/custom}{formatting}
\import{\ProjectRootPath/preambles/custom}{list_setup}

% TODO: перевести команды на современный стиль объявления

\begin{document}
    $x \parder{z} {x} + y \parder{z} {y} = z - xy$.
    
    $\dvfrac{\di x} {x} = \dvfrac{\di y} {y} = \dvfrac{\di z} {z - xy}$.
    
    \begin{enumerate}
        \item $\dvfrac{\di x} {x} = \dvfrac{\di y} {y}$.
        $\dvfrac{y}{x} = A$.
        \item $\dvfrac{\di x} {x} = \dvfrac{\di z} {z - xy}$.
        $xz\` = z - xy$. Это линейное уравнение, если учесть $y = Ax$.
        Подставив, получим
        $\dvfrac {xz\` - z} {\pow{x}{2}} = -A$.
        $\dvfrac {z}{x} = -Ax + B$.
        $\dvfrac{z + xy}{x} = B$.
    \end{enumerate}

    Результатом являются два интеграла 
    $\tsys{
        y = Ax \\
        z + xy = Bx,
    }$
    куда подставим $x = 2$ и $z = \pow{y}{2} + 1$.
    $\tsys{
        y = 2A \\
        \pow{y}{2} + 1 + 2y = 2B,
    }$ поэтому зависимость между интегралами такова: $\br{2A + 1}^2 = 2B$.
    В итоге $\br{2 \cdot \dvfrac {y}{x} + 1}^2 = 2 \cdot \dvfrac {z + xy}{x}$;
    $\br{2y + x}^2 = 2x\br{z + xy}$.
    В явном виде имеем функцию
    $z\br{x; y} = \dvfrac {\br{2y + x}^2}{2x} - xy$.
\end{document}
